% Statistiques -- Ajustement linéaire (méthode de Mayer, moindres carrés)
% Transcription LaTeX des notes manuscrites (stats.pdf)
\documentclass[11pt,a4paper]{article}

\usepackage[T1]{fontenc}
\usepackage[utf8]{inputenc}
\usepackage[french]{babel}
\usepackage{lmodern}
\usepackage{microtype}
\usepackage[a4paper,margin=2.5cm]{geometry}
\usepackage{amsmath,amssymb}
\usepackage{array}
\usepackage{enumitem}
\usepackage{needspace}

\frenchsetup{StandardLists=true}
\setlist[itemize]{label=\textbullet}

% Numérotation : I, II pour les parties ; 1, 2... pour les étapes
\renewcommand{\thesection}{\Roman{section}}
\renewcommand{\thesubsection}{\arabic{subsection}}

\newcommand{\abs}[1]{\left|#1\right|}

% Points à compléter
\newcommand{\blanc}[1][1.8cm]{\makebox[#1]{\dotfill}}

% Rubriques de l'activité
\newcommand{\rubrique}[1]{\par\medskip\noindent\textbf{#1}\par\smallskip}

% Tableaux
\newcolumntype{C}[1]{>{\centering\arraybackslash}p{#1}}
\renewcommand{\arraystretch}{1.3}
\newcommand{\barre}{---}

\setlength{\parindent}{0pt}
\setlength{\parskip}{0.45em}

\begin{document}

\begin{center}
  {\LARGE\bfseries Statistiques}
\end{center}

% =================================================================
\section{Ajustement linéaire par la méthode de Mayer}
% =================================================================

On a recueilli les données concernant les notes $x$ de mathématiques et $y$ de
physique d'un certain nombre d'élèves de la classe. Les données sont présentées
dans un tableau :

\begin{center}
  \small\setlength{\tabcolsep}{3pt}
  \begin{tabular}{|c|*{15}{C{0.62cm}|}}
    \hline
    $x$ & 10 & 12 & 10 & 13 & 09 & 11 & 12 & 11 & 13 & 12 & 14 & 10 & 15 & 10 & 12\\ \hline
    $y$ & 11 & 13 & 10 & 13 & 10 & 10 & 12 & 12 & 14 & 13 & 15 & 10 & 16 & 11 & 13\\ \hline
  \end{tabular}
\end{center}

\subsection{Situation problème}

Peut-on prédire la note $y$ de physique d'un élève qui a obtenu une note $x$ en
mathématiques ? En d'autres termes, peut-on trouver $a$ et $b$ réels tels que
\[
  y = ax + b \;?
\]

\subsection{Activité d'apprentissage}

\rubrique{Série marginale en $x$}

\begin{center}
  \begin{tabular}{|c|*{7}{C{0.8cm}|}c|}
    \hline
    $x_i$ & 09 & 10 & 11 & 12 & 13 & 14 & 15 & Total\\ \hline
    $n_i$ & & & & & & & & 15\\ \hline
  \end{tabular}
\end{center}
La moyenne $\bar{x}$ en mathématiques est
$\bar{x} = \blanc[3cm] = \blanc$.

\rubrique{Série marginale en $y$}

\begin{center}
  \begin{tabular}{|c|*{7}{C{0.8cm}|}c|}
    \hline
    $y_i$ & 10 & 11 & 12 & 13 & 14 & 15 & 16 & Total\\ \hline
    $n_i$ & & & & & & & & 15\\ \hline
  \end{tabular}
\end{center}
La moyenne $\bar{y}$ en physique est
$\bar{y} = \blanc[3cm] = \blanc$.

\rubrique{Coordonnées du point moyen $G$}

Le point moyen est $G(\bar{x}\,;\,\bar{y})$, donc
$G\bigl(\blanc\,;\,\blanc\bigr)$.

\rubrique{Nuage de points}

C'est la représentation des points $(x_i\,;\,y_i)$ du tableau à double entrée
des données.

Représenter les points.


Ces points semblent être regroupés autour d'une droite affine ; donc
l'ajustement linéaire de $y$ en $x$ est possible. En posant
$(\Delta) : y = ax + b$ cette droite, le but est de rechercher $a$ et $b$.

\rubrique{Détermination de la droite d'ajustement linéaire de $y$ en $x$ par la
méthode de Mayer}

Divisons la série statistique en deux sous-séries de même effectif ou
d'effectifs différents de 1 (ici 8 et 7).

\begin{center}
  \small
  \setlength{\tabcolsep}{4pt}
  \begin{tabular}{|c|*{8}{C{0.55cm}|}}
    \hline
    $x_i$ & 10 & 12 & 10 & 13 & 09 & 11 & 12 & 11\\ \hline
    $y_i$ & 11 & 13 & 10 & 13 & 10 & 10 & 12 & 12\\ \hline
  \end{tabular}
  \hfill
  \begin{tabular}{|c|*{7}{C{0.55cm}|}}
    \hline
    $x_i$ & 13 & 12 & 14 & 10 & 15 & 10 & 12\\ \hline
    $y_i$ & 14 & 13 & 15 & 10 & 16 & 11 & 13\\ \hline
  \end{tabular}
\end{center}

Pour chaque sous-série, calculons les points moyens $G_1$ et $G_2$ :
\begin{align*}
  G_1(\bar{x}_1\,;\,\bar{y}_1) &\quad\text{donc}\quad G_1\bigl(\blanc\,;\,\blanc\bigr)\\
  G_2(\bar{x}_2\,;\,\bar{y}_2) &\quad\text{donc}\quad G_2\bigl(\blanc\,;\,\blanc\bigr).
\end{align*}
\begin{enumerate}[label=\roman*)]
  \item Calculer
        $a = \dfrac{\bar{y}_1 - \bar{y}_2}{\bar{x}_1 - \bar{x}_2} = \blanc[2.4cm] = \blanc$.
  \item Calculer $b = \bar{y}_1 - a\bar{x}_1$ ou $b = \bar{y}_2 - a\bar{x}_2$ ou
        $b = \bar{y} - a\bar{x}$.
  \item L'équation de la droite d'ajustement linéaire par la méthode de Mayer
        est $y = ax + b$.
\end{enumerate}

\textbf{Remarque :} c'est en fait la droite $(G_1G_2)$ ; elle passe toujours par
le point $G(\bar{x}\,;\,\bar{y})$.

\subsection{Retenons}

\begin{itemize}
    \item Lorsque le nuage de points a la \og forme\fg{} d'une droite, on peut
          faire un ajustement linéaire de $y$ en $x$.
    \item La méthode de Mayer consiste à diviser la série statistique en deux
          sous-séries, à déterminer les coordonnées des deux points moyens
          $G_1$ et $G_2$, puis à déterminer l'équation de la droite
          $(G_1G_2) : y = ax + b$ où
          \[
            a = \frac{\bar{y}_1 - \bar{y}_2}{\bar{x}_1 - \bar{x}_2}
            \qquad\text{et}\qquad
            b = \bar{y}_1 - a\bar{x}_1 .
          \]
\end{itemize}

\needspace{10\baselineskip}
\subsection{Exercice d'application}

  \begin{itemize}
    \item Pour l'activité d'apprentissage, l'équation de la droite
          d'ajustement linéaire est ${y = ax + b} = \blanc[4cm]$.
    \item Déterminer la note théorique de physique d'un élève qui a obtenu 17
          en mathématiques.
    \item Déterminer la note théorique de mathématiques d'une élève qui a
          obtenu 8 en physique.
    \item Vérifier que $G$ appartient à cette droite.
  \end{itemize}

% =================================================================
\section{Ajustement linéaire par la méthode des moindres carrés}
% =================================================================

À partir du tableau des données :

\begin{center}
  \small\setlength{\tabcolsep}{3pt}
  \begin{tabular}{|c|*{15}{C{0.62cm}|}}
    \hline
    $x_i$ & 10 & 12 & 10 & 13 & 09 & 11 & 12 & 11 & 13 & 12 & 14 & 10 & 15 & 10 & 12\\ \hline
    $y_i$ & 11 & 13 & 10 & 13 & 10 & 10 & 12 & 12 & 14 & 13 & 15 & 10 & 16 & 11 & 13\\ \hline
  \end{tabular}
\end{center}

\subsection{Activité d'apprentissage}

$\bullet$ Calculer la moyenne
\[
  \bar{x} = \frac{\sum_{i=1}^{15} x_i}{N} = \frac{\blanc}{15} = \blanc
\]
et la moyenne
\[
  \bar{y} = \frac{\sum_{i=1}^{15} y_i}{N} = \frac{\blanc}{15} = \blanc\,;
\]
c'est la méthode directe lorsqu'on a un tableau à deux lignes ; mais on peut
aussi utiliser les séries marginales en $X$ et en $Y$.

$\bullet$ Calculer la covariance de $(X,Y)$ en complétant le tableau et en
appliquant les formules.

\begin{center}
  \small\setlength{\tabcolsep}{3pt}
  \begin{tabular}{|c|*{15}{C{0.6cm}|}C{1.1cm}|}
    \cline{17-17}
    \multicolumn{16}{c|}{} & Total\\ \hline
    $x_i$ & 10 & 12 & 10 & 13 & 09 & 11 & 12 & 11 & 13 & 12 & 14 & 10 & 15 & 10 & 12 & \barre\\ \hline
    $y_i$ & 11 & 13 & 10 & 13 & 10 & 10 & 12 & 12 & 14 & 13 & 15 & 10 & 16 & 11 & 13 & \barre\\ \hline
    $x_i \times y_i$ & & & & & & & & & & & & & & & & \\ \hline
  \end{tabular}
\end{center}

$\bullet$ La covariance de $(X,Y)$ est
\[
  \operatorname{cov}(X,Y) = \frac{\sum_{i=1}^{15} x_i y_i}{N} - \bar{x}\,\bar{y}
  = \frac{\blanc}{15} - \bigl(\blanc\bigr)\bigl(\blanc\bigr)
  = \blanc .
\]

$\bullet$ Calcul de la variance de $X$, $V(X)$ ; de l'écart type de
$X$, $\sigma_X$ ; de la variance de $Y$, $V(Y)$ ; de l'écart type de $Y$,
$\sigma_Y$.

Ici, on utilise la série marginale en $X$ et la série marginale en $Y$.

\begin{center}
  \begin{tabular}{|c|*{7}{C{0.85cm}|}C{1.1cm}|}
    \hline
    $x_i$ & 09 & 10 & 11 & 12 & 13 & 14 & 15 & Total\\ \hline
    $n_i$ & 1 & 4 & 2 & 4 & 2 & 1 & 1 & 15\\ \hline
    $n_i \times x_i$ & & & & & & & & \\ \hline
    $x_i^2$ & & & & & & & & \barre\\ \hline
    $n_i \times x_i^2$ & & & & & & & & \\ \hline
  \end{tabular}

  \bigskip
  \begin{tabular}{|c|*{7}{C{0.85cm}|}C{1.1cm}|}
    \hline
    $y_i$ & 10 & 11 & 12 & 13 & 14 & 15 & 16 & Total\\ \hline
    $n_i$ & 4 & 2 & 2 & 4 & 1 & 1 & 1 & 15\\ \hline
    $n_i \times y_i$ & & & & & & & & \\ \hline
    $y_i^2$ & & & & & & & & \barre\\ \hline
    $n_i \times y_i^2$ & & & & & & & & \\ \hline
  \end{tabular}
\end{center}

\needspace{14\baselineskip}
On a alors :

\noindent
\begin{minipage}[t]{0.48\linewidth}
  La moyenne de $X$ est
  \[
    \bar{x} = \frac{\sum_{i=1}^{7} n_i x_i}{N} = \frac{\blanc}{15} = \blanc
  \]
  La variance de $X$ est
  \[
    V(X) = \frac{\sum_{i=1}^{7} n_i x_i^2}{N} - (\bar{x})^2 = \blanc
  \]
  L'écart type de $X$ est
  \[
    \sigma_X = \sqrt{V(X)} = \blanc = \blanc
  \]
\end{minipage}%
\hfill
\begin{minipage}[t]{0.48\linewidth}
  La moyenne de $Y$ est
  \[
    \bar{y} = \frac{\sum_{i=1}^{7} n_i y_i}{N} = \frac{\blanc}{15} = \blanc
  \]
  La variance de $Y$ est
  \[
    V(Y) = \frac{\sum_{i=1}^{7} n_i y_i^2}{N} - (\bar{y})^2 = \blanc
  \]
  L'écart type de $Y$ est
  \[
    \sigma_Y = \sqrt{V(Y)} = \blanc = \blanc
  \]
\end{minipage}

\medskip
\textbf{Exercice.} Avec les données précédentes, calculer
  \[
    a = \frac{\operatorname{cov}(X,Y)}{V(X)} = \blanc[2.4cm]
    \qquad;\qquad
    a' = \frac{\operatorname{cov}(X,Y)}{V(Y)} = \blanc[2.4cm]
  \]
  \begin{itemize}
    \item Déterminer le point moyen $G(\bar{x}\,;\,\bar{y})$ :
          $G\bigl(\blanc\,;\,\blanc\bigr)$.
    \item Calculer $a \times a' = \blanc[2.4cm] = \blanc$.
    \item Calculer
          $r = \dfrac{\operatorname{cov}(X,Y)}{\sqrt{V(X)} \times \sqrt{V(Y)}}$.
  \end{itemize}

\subsection{\texorpdfstring{Corrélation linéaire entre $X$ et $Y$}{Corrélation linéaire entre X et Y}}

\textbf{Retenons.}
\begin{itemize}
    \item Le coefficient de corrélation linéaire entre $X$ et $Y$ est
          \[
            r = \frac{\operatorname{cov}(X,Y)}{\sqrt{V(X)} \times \sqrt{V(Y)}} .
          \]
    \item Si $\abs{r} \le 0{,}5$, on dit que la corrélation est faible.

          Si $0{,}7 \le \abs{r} < 1$, on dit que la corrélation est bonne ou
          forte.

          Si $\abs{r} = 1$, on dit que la corrélation est parfaite.
    \item L'ajustement linéaire se justifie si $0{,}7 \le \abs{r} \le 1$.
\end{itemize}

\subsection{\texorpdfstring{Droite de régression de $y$ en $x$ ou droite
d'ajustement linéaire obtenue par la méthode des moindres carrés}{Droite de
régression de y en x ou droite d'ajustement linéaire obtenue par la méthode des
moindres carrés}}

On calcule
\[
  a = \frac{\operatorname{cov}(X,Y)}{V(X)}
  \qquad\text{et}\qquad
  G(\bar{x}\,;\,\bar{y}),
\]
alors
\[
  (\Delta) : y - \bar{y} = a(x - \bar{x}) ;
\]
on en déduit $y = ax + \bar{y} - a\bar{x}$, ou bien
\[
  y = ax + b
  \qquad\text{avec}\qquad
  a = \frac{\operatorname{cov}(X,Y)}{V(X)}
  \quad\text{et}\quad
  b = \bar{y} - a\bar{x} .
\]

\textbf{Remarque :} pour la droite de régression de $x$ en $y$, on calcule
\[
  a' = \frac{\operatorname{cov}(X,Y)}{V(Y)}
\]
et
\[
  (\Delta') : x - \bar{x} = a'(y - \bar{y}),
  \qquad\text{c'est-à-dire}\qquad
  x = a'y + \bar{x} - a'\bar{y} .
\]

\textbf{Attention :}
\begin{itemize}
  \item le coefficient directeur de $(\Delta')$ est
        $\dfrac{1}{a'} = \dfrac{V(Y)}{\operatorname{cov}(X,Y)}$ ;
  \item on a :
        $a \times a' = \dfrac{\bigl(\operatorname{cov}(X,Y)\bigr)^2}{V(X)\,V(Y)} = r^2$ ;
  \item justifier que $a$, $a'$ et $\operatorname{cov}(X,Y)$ ont le même signe.
\end{itemize}

\subsection{Exercice d'application}

  \begin{enumerate}[label=\alph*)]
    \item Ayant calculé le coefficient de corrélation linéaire de l'activité
          d'apprentissage, qualifier la corrélation entre $X$ et $Y$.
          \begin{enumerate}[label=\roman*)]
            \item Écrire l'équation de la droite de régression de $y$ en $x$.
            \item Un élève a 16 en mathématiques ; estimer sa note théorique de
                  physique.
            \item Une élève a 8 en physique ; estimer sa note théorique de
                  mathématiques.
          \end{enumerate}

    \item $(\Delta_1) : y = 6x - 38$ et $(\Delta_2) : x = 0{,}135y + 6{,}65$
          sont les droites de régression d'une série double. Déterminer $G$ et
          $r$. Que dire de la corrélation de $X$ et $Y$ ?
          \begin{itemize}
            \item Même exercice pour $(\Delta_1) : y = 0{,}8x - 3$ et
                  $(\Delta_2) : x = 0{,}9y + 5{,}5$.
            \item Même exercice pour $(\Delta_1) : y = 1{,}5x + 3$ et
                  $(\Delta_2) : x = 0{,}54y - 0{,}86$.
          \end{itemize}

    \item La droite de régression de $y$ en $x$ d'équation
          $y = \dfrac{27}{14}x + \dfrac{16}{7}$ a été obtenue à partir du
          tableau ci-dessous ; les réels $a$ et $b$ ont été malheureusement
          effacés. Déterminer les nombres $a$ et $b$.
          \begin{center}
            \begin{tabular}{|c|*{7}{C{0.8cm}|}}
              \hline
              $x_i$ & 1 & 2 & 3 & 4 & 5 & 6 & 7\\ \hline
              $y_i$ & 4 & 6 & $a$ & 10 & 13 & $b$ & 15\\ \hline
            \end{tabular}
          \end{center}

    \needspace{4\baselineskip}
    \item On considère le tableau suivant :
          \begin{center}
            \begin{tabular}{|c|*{7}{C{1cm}|}}
              \hline
              $x_i$ & 1 & 2 & 3 & 4 & 5 & 6 & 7\\ \hline
              $y_i$ & $a$ & $a+2$ & $2a$ & $b-4$ & $b-1$ & $b$ & $b+1$\\ \hline
            \end{tabular}
          \end{center}
          La droite de régression de $y$ en $x$ a pour équation
          $27x - 14y + 32 = 0$.
          \begin{enumerate}[label=\roman*)]
            \item Calculer $a$ et $b$.
            \item Restituer le tableau original.
          \end{enumerate}

    \item On considère le tableau suivant :
          \begin{center}
            \begin{tabular}{|c|*{7}{C{0.8cm}|}}
              \hline
              $x_i$ & 5 & 10 & 11 & 8 & 7 & 6 & 9\\ \hline
              $y_i$ & 11 & 21 & 23 & 17 & 15 & 13 & 19\\ \hline
            \end{tabular}
          \end{center}
          \begin{enumerate}[label=\roman*)]
            \item La corrélation entre $x$ et $y$ est-elle parfaite ?
            \item Écrire l'équation de la droite de régression de $y$ en $x$.
          \end{enumerate}
  \end{enumerate}

\end{document}
